% Inactive consequence of the retained curvature-signature exhaustion.
\subsection{External tensor products and complement convexity}
\label{subsec:external-tensor-complement-curvature}

On a product base, separate positive curvature traces give separate positive tangent directions. The external tensor metric has no curvature terms between different factors, so this count applies to every fibre vector, including those that are not decomposable tensors.

Let $k\geq2$. For $1\leq i\leq k$, let $X_i$ be a connected compact complex manifold of dimension $n_i\geq1$, let $E_i\to X_i$ have rank $r_i\geq1$, and choose a smooth Hermitian metric $h_i$ and a Hermitian form $\omega_i$. With $p_i:X=\prod_iX_i\to X_i$ the projections, put
\[
 E=\bigotimes_{i=1}^kp_i^*E_i,\qquad h=\bigotimes_{i=1}^kp_i^*h_i,
 \qquad n=\sum_i n_i,\quad r=\prod_i r_i.
\]
For an endomorphism $A_i$ of $E_i$, write $\widehat A_i$ for its action on the $i$th tensor factor, with the identity on every other factor. Use the Chern and mean curvatures of \cref{subsec:projective-complement-curvature-index}.

\begin{lemma}\label{lem:external-tensor-curvature-traces}
The curvature of $h$ is
\[
 F_h=\sum_i\widehat{p_i^*F_{h_i}}.
\]
For every $v\in E_x$, the coefficient form of $\alpha_{h,v}$ is block diagonal for $T_x^{1,0}X=\bigoplus_iT_{x_i}^{1,0}X_i$, and
\[
 \operatorname{tr}_{\omega_i}\bigl(\alpha_{h,v}|_{T_{x_i}X_i}\bigr)
 =h\bigl(\widehat K_i v,v\bigr),\qquad K_i=K_{\omega_i,h_i}.
\]
\end{lemma}
\begin{proof}
In product holomorphic frames the metric matrix is $M=\bigotimes_iM_i$. Its connection matrix is
\[
 M^{-1}\partial M=\sum_i\widehat{M_i^{-1}\partial M_i}.
\]
Each summand depends only on the coordinates of $X_i$. Applying $\bar\partial$ gives the curvature formula and shows that all mixed factor components vanish. Pairing with $v$ gives the block diagonal form. Contracting its $i$th block with $\omega_i$ commutes with the identity actions on the remaining tensor factors and gives the trace identity.
\end{proof}

\begin{theorem}\label{thm:external-tensor-complement-convexity}
If every $K_i$ is positive definite, then
\[
 \mathbb P(\mathcal O_X\oplus E)\setminus\mathbb P(\mathcal O_X)
\]
is $(n+r-k)$-convex.
\end{theorem}
\begin{proof}
At each point choose an $h_i$-orthonormal eigenbasis of $K_i$. The resulting product basis is $h$-orthonormal, and $\widehat K_i$ has the positive eigenvalues of $K_i$, repeated for the other tensor factors. Hence $h(\widehat K_i v,v)>0$ for every nonzero $v$, without a decomposability assumption.

By \cref{lem:external-tensor-curvature-traces}, each tangent block of $\alpha_{h,v}$ has positive trace and thus a positive eigenvalue. The block diagonal form has at least $k$ positive eigenvalues. Since $k\leq n$, apply \cref{thm:bundle-curvature-complement-convexity} with $s=k$.
\end{proof}

\begin{corollary}\label{cor:stable-factor-complement-convexity}
Suppose that each $\omega_i$ is K\"ahler, that $E_i$ is slope-stable with respect to $\omega_i$, and that
\[
 \int_{X_i}c_1(E_i)\wedge\omega_i^{n_i-1}>0.
\]
Then the external tensor complement is $(n+r-k)$-convex. For $k$ stable ample Hirzebruch examples of \cref{cor:hirzebruch-exact-complement-index}, with their specified polarizations, this gives $(2^k+k)$-convexity.
\end{corollary}
\begin{proof}
As in the proof of \cref{cor:positive-mean-complement-convexity}, choose the Hermitian--Einstein metric on each $E_i$. Its mean curvature is $\lambda_i\operatorname{Id}_{E_i}$, with $\lambda_i>0$ by the positive degree and integrated Chern trace identity. Apply \cref{thm:external-tensor-complement-convexity}. For the Hirzebruch factors, $n_i=r_i=2$, so $n+r-k=2k+2^k-k=2^k+k$.
\end{proof}
